\section{Why PEM Works}
\label{sec:theory}

This section gives a formal account of why Post-hoc Entropy Minimization is well-conditioned in the regime of a converged SSL checkpoint and ill-conditioned elsewhere. The argument has three parts: a gradient computation showing that the binary-entropy gradient is an implicit boundary-shell gate (Sec.~\ref{sec:theory:p1p}); a connection between the entropy gradient and the classical cluster assumption specialized to 3D segmentation at the organ surface (Sec.~\ref{sec:theory:cluster}); and a comparison with training-time entropy minimization and pseudo-labeling that explains the tradeoffs we observe in the experiments (Sec.~\ref{sec:theory:comparison}). All three parts predict both failure modes we characterize in Section~\ref{sec:discussion}, and both hold on the specific checkpoints we evaluate in Section~\ref{sec:experiments}.

\subsection{The $p(1{-}p)$ factor as an implicit boundary-shell gate}
\label{sec:theory:p1p}

Consider the simplest nontrivial case: a binary voxel classifier with logit $z \in \mathbb{R}$ and probability $p\!=\!\sigma(z)\!=\!(1 + e^{-z})^{-1}$. The per-voxel Shannon entropy of this distribution is
\begin{equation}
H(p) \;=\; -p \log p \;-\; (1-p) \log(1-p).
\end{equation}
A direct computation gives $\partial H/\partial p \;=\; -\log p + \log(1-p) \;=\; \log\!\big((1-p)/p\big) \;=\; -z$, and therefore by the chain rule
\begin{equation}
\frac{\partial H(p(z))}{\partial z} \;=\; \frac{\partial H}{\partial p}\cdot\frac{\partial p}{\partial z} \;=\; (-z)\cdot p(1-p) \;=\; -\,z\,p(1-p).
\label{eq:dHdz}
\end{equation}
Equation~\ref{eq:dHdz} is the central object of this section. It expresses the per-voxel entropy-minimization gradient with respect to the logit as the product of two terms: the logit itself, which is large in magnitude on voxels where the network is confident, and the variance factor $p(1{-}p)$, which is large only on voxels where the network is undecided. The product of these two terms peaks at intermediate logit values --- specifically at $|z| \approx 1.543$, which corresponds to $p(1{-}p) \approx 0.209$ and $|{\partial H/\partial z}| \approx 0.322$ --- and decays to zero at both extremes. The peak is broad: the gradient retains at least half its maximum for $p$ between roughly $0.12$ and $0.88$, and drops below $10\%$ of its maximum only for $p$ outside $[0.02, 0.98]$.

\begin{proposition}[Implicit boundary-shell gating]
\label{prop:gating}
Let $f_{\theta^*}$ be a binary segmentation network with voxel logits $z_i$ and softmax $p_i\!=\!\sigma(z_i)$, and let $\mathcal{L}_{\mathrm{PEM}}$ be the PEM loss (Eq.~\ref{eq:pem}) with the mask $m_i$ removed (i.e.\ $\tau\!=\!1$). The magnitude of the per-voxel gradient contribution to $\partial \mathcal{L}_{\mathrm{PEM}}/\partial \theta$ is upper-bounded by a quantity proportional to $|z_i|\cdot p_i(1{-}p_i)$, which vanishes as $p_i\!\to\!0$ or $p_i\!\to\!1$ and is therefore dominated by the contribution from the uncertain shell $\{i : p_i\!\approx\!0.5\}$. In particular, no explicit confidence mask is required for the loss to be localized at the decision surface; the $p(1{-}p)$ factor provides this localization implicitly.
\end{proposition}

Proposition~\ref{prop:gating} has three consequences that explain the observed behavior of PEM without reference to any explicit masking.

\paragraph{Consequence 1: no class-balance regularizer is needed.}
A classical objection to unmasked entropy minimization as an SSL loss is that it admits a trivial minimizer at $p_i\!=\!\mathbb{1}[\hat c_i]$ for any choice of class $\hat c_i$, so that a network minimizing entropy with no other constraint will collapse to predicting a single class everywhere~\cite{grandvalet2004entropy}. Training-time entropy minimization circumvents this by coupling the entropy term to a supervised cross-entropy that disagrees with the collapsed prediction, and by adding a class-balance regularizer that explicitly penalizes deviations from the expected class prior. Neither stabilizer is needed in the post-hoc regime, because the $p(1{-}p)$ factor already vanishes on saturated voxels: a voxel that has already been pushed to $p\!\approx\!0$ or $p\!\approx\!1$ by the base SSL training cannot be pushed further by the PEM gradient, and the optimization therefore cannot collapse a converged checkpoint into a single-class predictor. The only voxels at which the gradient is substantial are those for which the base model was undecided to begin with, and the only direction it can push them is toward the class currently assigned slightly higher probability.

\paragraph{Consequence 2: the soft gradient shrinks as the base model saturates.}
As the base checkpoint becomes more confident, the mean of $p(1{-}p)$ over the volume shrinks toward zero, and with it the magnitude of the PEM gradient. Concretely, let $\bar V \!=\! \mathbb{E}_i[p_i(1{-}p_i)]$ be the mean Bernoulli variance over voxels; then $\|\partial \mathcal{L}_{\mathrm{PEM}} / \partial \theta\|$ scales at most as $O(\bar V \cdot \|\nabla_\theta z\|_\infty)$. A fully saturated network has $\bar V\!=\!0$ and therefore $\|\partial \mathcal{L}_{\mathrm{PEM}} / \partial \theta\|\!=\!0$: it cannot be improved by PEM at all. A partially saturated network has $\bar V\!>\!0$ but small, and its PEM improvement is proportional to this small $\bar V$. This is the formal content of the second failure mode we observe in Section~\ref{sec:discussion}: as $\rho_f(0.99)\!\to\!1$ and $\bar V\!\to\!0$, PEM gains diminish, and pseudo-label fine-tuning (whose gradients are $O(1)$ regardless of $\bar V$) becomes competitive. The LA 10\% result is exactly this regime.

\paragraph{Consequence 3: the gradient direction is aligned with logit descent on boundary voxels.}
On a boundary voxel with $p_i\!\approx\!0.5$, the gradient $-z_i\,p_i(1{-}p_i)$ has the sign of $z_i$ and therefore pushes $z_i$ further from zero in whichever direction it is currently trending. In other words, PEM commits the voxel to the class it is already slightly more likely to predict, with a step size proportional to the current undecidedness. If the base SSL training has placed the boundary in the right \emph{neighborhood} --- which the convergence of the SSL loss guarantees to first order --- then the PEM step refines the boundary rather than relocating it. If the base SSL training has placed the boundary in the wrong neighborhood (e.g.\ a supervised-only network with diffuse $\rho_f$), the PEM step reinforces the error; this is the first failure mode.

\subsection{Connection to the cluster assumption at the physical organ surface}
\label{sec:theory:cluster}

Proposition~\ref{prop:gating} is a statement about gradient mechanics; it does not, by itself, tell us \emph{why} decreasing entropy at boundary voxels should improve segmentation quality rather than merely make the network more confident about whatever its current predictions are. The connection between entropy minimization and actual accuracy improvement comes from the classical cluster assumption of Chapelle et al.~\cite{chapelle2006ssl}, which states that classes form dense clusters in feature space and that the optimal decision boundary lies in the low-density regions between them. Grandvalet and Bengio~\cite{grandvalet2004entropy} observed that minimizing the conditional entropy $H(Y \mid X)$ on unlabeled data implements this principle: a classifier that maximizes its margin around the current decision surface will tend to push that surface into low-density regions of feature space. We extend this argument in two directions specific to 3D medical image segmentation.

\paragraph{Feature-space boundaries and anatomical boundaries coincide.}
In 3D medical image segmentation, the network's feature-space decision surface is not arbitrary: it is spatially co-located with the physical organ surface. The reason is structural rather than learned. Boundary voxels, in the physical sense of voxels lying on or near the interface between two tissues, are precisely the voxels where the per-voxel feature distribution is bimodal. Partial-volume averaging causes the HU or intensity signal at such a voxel to be a mixture of the two neighboring tissues; scanner point spread functions blur the signal further; soft-tissue contrast on MR and CT is often weak precisely at the clinically relevant interfaces; and biological variability causes the precise location of the surface to differ across volumes by a few voxels even in the absence of pathology. A deep network trained on this data learns that these voxels are uncertain because their input features are bimodal in its feature space --- not because the network is ignorant of them, but because the information present in the input is genuinely ambiguous. The network's feature-space decision surface therefore lies \emph{inside} the physically bimodal voxels, which are the same voxels as the anatomical boundary.

\paragraph{Entropy minimization at a physically meaningful surface is boundary refinement.}
Under this alignment, the gradient step implied by Eq.~\ref{eq:dHdz} has a direct anatomical interpretation. For each boundary voxel, the sign of the logit $z_i$ tells us which class is currently more probable given the voxel's features; the $p(1{-}p)$ factor tells us how undecided the network currently is; and the product of the two, taken as a gradient step, moves the voxel's logit in the direction of the currently more-probable class by an amount proportional to the current undecidedness. Anatomically, this is equivalent to sliding the predicted organ surface toward the locally higher-probability tissue class by an amount proportional to the local ambiguity. For a well-converged SSL network, the sign of $z_i$ is correct in expectation at most boundary voxels (otherwise the network would not have converged), and the PEM step therefore slides the predicted surface toward the ground-truth surface. For a supervised-only network whose boundary has been placed on the wrong side of some anatomically ambiguous voxels, the sign of $z_i$ is incorrect for those voxels, and the PEM step slides them further from the truth. The distinction between these two cases is \emph{predicted} by $\rho_f(0.99)$: a high $\rho_f$ means the uncertain voxels are confined to a thin physical shell, which by the argument above is the same shell as the anatomical organ surface; a low $\rho_f$ means the uncertain voxels are scattered across interior regions whose logits do not correspond to any anatomical surface, and for which the sign of $z_i$ is essentially random.

\paragraph{Why this is specifically a 3D segmentation argument.}
The coincidence between feature-space and physical boundaries is not a general property of classification tasks. In image-level classification, for instance, the feature-space decision surface has no natural spatial interpretation, and entropy minimization is a margin-maximization argument that lives entirely in feature space. In natural-image semantic segmentation, the coincidence is weaker because semantic boundaries (object outlines) can be drawn sharply by human annotators, so the feature-space boundary has more freedom to deviate from any particular pixel. In 3D medical segmentation, the coincidence is essentially perfect: the anatomical surface is defined by the imaging physics as the set of voxels where the measurement itself is ambiguous, and no other set of voxels can plausibly be placed on the decision surface by a well-converged network. This is why PEM is particularly effective on medical 3D volumes and why the $\rho_f$ diagnostic is a reliable applicability check in this setting.

\subsection{Why not train-time EM, and why not pseudo-labels?}
\label{sec:theory:comparison}

The above argument explains why post-hoc entropy minimization is well-conditioned in the SSL convergence regime. It does not yet explain why we cannot achieve the same effect by adding an entropy term to the SSL training loss in the first place, nor why we should prefer a soft entropy gradient to the hard pseudo-label gradient of Lee~\cite{lee2013pseudolabel}. We address both questions in turn.

\paragraph{Training-time entropy minimization is a balance problem.}
When an entropy term is added to a supervised loss during training, its strength must be carefully balanced against the cross-entropy. Too weak, and it has no effect; too strong, and it drives the posterior toward degenerate collapse before the supervised loss has had a chance to place the decision boundary anywhere anatomically meaningful. The literature is full of devices for navigating this balance: warmup schedules that gradually ramp the entropy coefficient from zero; confidence thresholds that only apply the loss to voxels above some margin; class-balance regularizers that penalize posterior skew; and curriculum learning over labeled fractions. Each device works, and each introduces hyperparameters. The post-hoc regime sidesteps this problem entirely: by the time PEM is applied, the supervised-coupled SSL training has already placed the decision boundary in a near-correct neighborhood, and the entropy loss is no longer competing with a supervised signal for optimization budget. It is operating on the residue, and the residue is in exactly the form that makes entropy minimization a refinement operator rather than a relocation operator. This is why PEM has no warmup, no curriculum, no class-balance term, and essentially one hyperparameter (the learning rate), whereas every training-time entropy method in the literature has several.

\paragraph{Pseudo-labeling is the hard-target limit.}
Pseudo-labeling~\cite{lee2013pseudolabel} replaces the soft entropy gradient with a hard target obtained by taking the argmax of the current prediction and training against it under cross-entropy. In the binary case with logit $z$ and probability $p$, the hard-target cross-entropy gradient is $\partial \mathrm{CE}(p, \mathbb{1}[z\!>\!0])/\partial z \!=\! p - \mathbb{1}[z\!>\!0]$, which is $O(1)$ in magnitude regardless of how close $p$ is to $0$ or $1$. Contrast this with the soft entropy gradient $-z\,p(1{-}p)$, which vanishes as $p\!\to\!0$ or $p\!\to\!1$. The pseudo-label gradient never saturates: it continues to push even fully committed voxels toward their current argmax, which can lead to over-commitment if the argmax was slightly wrong. The soft entropy gradient is conservative: it refines undecided voxels and leaves committed voxels alone, trading the risk of over-commitment for a shrinking signal in the saturation limit.

\paragraph{Predicted tradeoff between PEM and pseudo-label fine-tune.}
The consequence is an interpretable tradeoff. When the residual entropy budget is substantial --- that is, when $\rho_f\!<\!1$ and the mean Bernoulli variance $\bar V$ is non-negligible --- PEM's soft gradient is large and conservative, and it outperforms PL-FT because it refines without committing errors. When the residual entropy budget is small --- that is, when $\rho_f\!\to\!1$ and $\bar V\!\to\!0$ --- PEM's soft gradient shrinks, and PL-FT's $O(1)$ gradient catches up. At some point the two methods become equivalent in expected improvement, and beyond it PL-FT can overtake PEM. This is exactly the ordering we observe in Section~\ref{sec:experiments}: PEM wins on Pancreas-CT 20\% ($\rho_f\!=\!0.894$, BCP baseline Dice $82.89$) and LA 5\% ($\rho_f\!=\!0.948$, baseline $87.32$); PL-FT marginally wins on LA 10\% ($\rho_f\!=\!0.945$, baseline $89.40$, where the total residual entropy budget is smallest because the Dice is highest). The tradeoff is quantitatively consistent with the $p(1{-}p)$ factor of Eq.~\ref{eq:dHdz} and is predictable from $\rho_f$ and the base Dice alone, without running either method.

\paragraph{Summary.}
PEM works in the SSL convergence regime for three interlocking reasons. First, the binary-entropy gradient $-z\,p(1{-}p)$ is an implicit boundary-shell gate that vanishes on saturated voxels, which is what prevents collapse and what makes an explicit mask unnecessary. Second, the convergence of the base SSL training concentrates residual entropy precisely at the physical organ surface (as measured by $\rho_f$), which is the set of voxels where entropy minimization's margin-maximization interpretation translates into anatomical boundary refinement. Third, a converged checkpoint is not in a balance-regime problem anymore: the supervised loss has already done its work, and the entropy loss can operate on the residue without needing warmups, schedules, or class-balance terms. The regime is uniquely favorable, and it is the regime that every publicly released SSL segmentation checkpoint on canonical 3D medical benchmarks appears to inhabit.
